\documentclass[a0paper,portrait]{article}
\usepackage[utf8]{inputenc}
\usepackage[T1]{fontenc}
\usepackage{lmodern}
\usepackage{amsmath,amssymb}
\usepackage{graphicx}
\usepackage{multicol}
\usepackage[margin=2.5cm]{geometry}
\usepackage{xcolor}
\usepackage{tikz}
\usepackage{enumitem}
\usepackage{booktabs}
\usepackage{microtype}

\pagestyle{empty}

% ──── Color Theme ────
\definecolor{headerblue}{HTML}{1B3A5C}
\definecolor{accentblue}{HTML}{2E86C1}
\definecolor{lightbg}{HTML}{EBF5FB}
\definecolor{darktext}{HTML}{2C3E50}
\definecolor{sectioncolor}{HTML}{1A5276}
\definecolor{boxbg}{HTML}{F4F8FB}

% ──── Section styling ────
\makeatletter
\renewcommand{\section}{\@startsection{section}{1}{0pt}%
  {-2ex plus -0.5ex minus -0.2ex}{1.5ex plus 0.3ex}%
  {\fontsize{40}{48}\selectfont\bfseries\color{sectioncolor}}}
\makeatother

\renewcommand{\familydefault}{\sfdefault}

% ──── Poster body font size ────
\newcommand{\posterbody}{\fontsize{28}{38}\selectfont}
\newcommand{\postertable}{\fontsize{26}{34}\selectfont}

\begin{document}

% ════════════════════════════════════════════
% HEADER
% ════════════════════════════════════════════
\begin{tikzpicture}[remember picture, overlay]
  \fill[headerblue] ([yshift=0cm]current page.north west) rectangle ([yshift=-18cm]current page.north east);
\end{tikzpicture}

\vspace*{-1cm}
\begin{center}
  {\color{white}\fontsize{80}{96}\selectfont\bfseries
  Machine Learning--Guided Drug Discovery:\\[0.4cm]
  Predicting Protein--Ligand Binding Affinity\\[0.4cm]
  with Graph Neural Networks}\\[2cm]
  {\color{white}\fontsize{40}{48}\selectfont
  \textbf{Elena Vasquez}$^1$, \textbf{Thomas Wright}$^{1,2}$, \textbf{Kenji Yamamoto}$^2$, \textbf{Prof. Laura Kingston}$^1$}\\[1cm]
  {\color{white!85}\fontsize{34}{42}\selectfont
  $^1$Department of Biochemistry, University of Oxford \quad\quad
  $^2$DeepMind, London, UK}\\[0.6cm]
  {\color{accentblue!40}\fontsize{30}{36}\selectfont
  Contact: elena.vasquez@bioch.ox.ac.uk \quad|\quad ICML 2025 -- Poster \#247}
\end{center}

\vspace{3cm}

% ════════════════════════════════════════════
% BODY -- Three Columns
% ════════════════════════════════════════════
\setlength{\columnsep}{3cm}
\begin{multicols}{3}

% ──── Introduction ────
\section*{Introduction}

\posterbody

Drug discovery is a lengthy and expensive process, with the average new drug requiring over 10 years and \$2.6 billion to develop. A critical bottleneck is the accurate prediction of protein--ligand binding affinity, which determines whether a candidate molecule will effectively interact with its target protein.

\vspace{1cm}

\textbf{Traditional approaches:}
\begin{itemize}[leftmargin=1.5em, itemsep=10pt]
  \item Molecular docking (AutoDock, Glide) -- fast but inaccurate
  \item Molecular dynamics simulations -- accurate but computationally prohibitive
  \item Classical QSAR models -- limited to predefined descriptors
\end{itemize}

\vspace{1cm}

\textbf{Our contribution:} We propose \textbf{AffinityGNN}, a graph neural network that operates directly on the 3D molecular graph of the protein--ligand complex to predict binding free energy ($\Delta G$) with near-experimental accuracy.

\vspace{1.5cm}

% ──── Methods ────
\section*{Methods}

\posterbody

\textbf{Graph Construction.} We represent the protein--ligand complex as a heterogeneous graph $G = (V_P \cup V_L, E)$ where:
\begin{itemize}[leftmargin=1.5em, itemsep=10pt]
  \item $V_P$: protein residue nodes (C$\alpha$ atoms)
  \item $V_L$: ligand heavy atom nodes
  \item $E$: edges based on spatial proximity ($< 8$\AA)
\end{itemize}

\vspace{1cm}

\textbf{Architecture.} AffinityGNN consists of:
\begin{enumerate}[leftmargin=1.5em, itemsep=10pt]
  \item Node feature encoder (atom type, charge, hybridization)
  \item 6 layers of message-passing with attention
  \item Graph-level readout with Set2Set pooling
  \item Prediction head: 3-layer MLP $\rightarrow \Delta G$
\end{enumerate}

\vspace{1cm}

\textbf{Message Passing:}
\begin{equation*}
  {\fontsize{32}{40}\selectfont h_i^{(l+1)} = \sigma\!\left( \sum_{j \in \mathcal{N}(i)} \alpha_{ij}^{(l)} W^{(l)} h_j^{(l)} + b^{(l)} \right)}
\end{equation*}

\vspace{0.5cm}
where $\alpha_{ij}$ are attention weights incorporating edge features (distance, angle).

\vspace{1cm}

\textbf{Training.} We use a combined loss:
\begin{equation*}
  {\fontsize{32}{40}\selectfont \mathcal{L} = \underbrace{\|\hat{y} - y\|_2^2}_{\text{MSE}} + \lambda \underbrace{(1 - \rho(\hat{y}, y))}_{\text{Correlation}}}
\end{equation*}

\vspace{1.5cm}

% ──── Data ────
\section*{Datasets}

\posterbody

\textbf{Training Data:}
\begin{itemize}[leftmargin=1.5em, itemsep=10pt]
  \item PDBbind v2020 refined set (5,316 complexes)
  \item BindingDB kinase subset (12,400 complexes)
  \item Custom curated GPCR dataset (3,200 complexes)
\end{itemize}

\vspace{1cm}

\textbf{Data Augmentation:}
\begin{itemize}[leftmargin=1.5em, itemsep=10pt]
  \item Random rotation and translation of coordinates
  \item Gaussian noise on atomic positions ($\sigma = 0.1$\AA)
  \item Subgraph sampling for large complexes
\end{itemize}

\columnbreak

% ──── Results ────
\section*{Results}

\posterbody

\textbf{Benchmark Performance} on PDBbind v2020 core set:

\vspace{1cm}

\begin{center}
\renewcommand{\arraystretch}{1.6}
{\postertable
\begin{tabular}{@{}lcc@{}}
\toprule
\textbf{Method} & \textbf{RMSE} & \textbf{Pearson $R$} \\
\midrule
AutoDock Vina  & 2.41 & 0.604 \\
RF-Score v3    & 1.82 & 0.713 \\
OnionNet-2     & 1.54 & 0.782 \\
PLIP-GNN       & 1.41 & 0.801 \\
\midrule
\textbf{AffinityGNN} & \textbf{1.18} & \textbf{0.862} \\
\bottomrule
\end{tabular}}
\end{center}

\vspace{2cm}

\textbf{Virtual Screening on DUD-E:}

\vspace{1cm}

\begin{center}
\renewcommand{\arraystretch}{1.6}
{\postertable
\begin{tabular}{@{}lcc@{}}
\toprule
\textbf{Target} & \textbf{AUC-ROC} & \textbf{EF$_{1\%}$} \\
\midrule
CDK2 (kinase)     & 0.94 & 42.3 \\
COX-2 (enzyme)    & 0.91 & 38.7 \\
EGFR (receptor)   & 0.93 & 45.1 \\
HIV-RT (viral)    & 0.89 & 31.4 \\
\midrule
\textbf{Average}  & \textbf{0.92} & \textbf{39.4} \\
\bottomrule
\end{tabular}}
\end{center}

\vspace{2cm}

\textbf{Experimental Validation.} We used AffinityGNN to screen 50,000 compounds against SARS-CoV-2 main protease (M$^{\text{pro}}$). The top 100 predictions were synthesized and tested:
\begin{itemize}[leftmargin=1.5em, itemsep=10pt]
  \item 23 compounds showed IC$_{50} < 10\,\mu$M
  \item 4 compounds showed IC$_{50} < 100\,$nM
  \item Best hit: IC$_{50} = 28\,$nM (comparable to nirmatrelvir)
\end{itemize}

\vspace{2cm}

\textbf{Generalization Across Targets:}

\vspace{1cm}

\begin{center}
\renewcommand{\arraystretch}{1.6}
{\postertable
\begin{tabular}{@{}lcc@{}}
\toprule
\textbf{Protein Family} & \textbf{RMSE} & \textbf{$N$} \\
\midrule
Kinases        & 1.12 & 847 \\
Proteases      & 1.21 & 523 \\
GPCRs          & 1.35 & 312 \\
Nuclear rec.   & 1.19 & 198 \\
\bottomrule
\end{tabular}}
\end{center}

\columnbreak

% ──── Analysis ────
\section*{Ablation \& Analysis}

\posterbody

\textbf{Component Contributions:}

\vspace{1cm}

\begin{center}
\renewcommand{\arraystretch}{1.6}
{\postertable
\begin{tabular}{@{}lc@{}}
\toprule
\textbf{Variant} & \textbf{RMSE} \\
\midrule
Full model                   & 1.18 \\
$-$ Attention mechanism      & 1.34 \\
$-$ Edge features            & 1.29 \\
$-$ Correlation loss         & 1.25 \\
$-$ Set2Set pooling          & 1.31 \\
GCN baseline (no attention)  & 1.52 \\
\bottomrule
\end{tabular}}
\end{center}

\vspace{2cm}

\textbf{Attention Visualization.} The learned attention weights highlight key interactions at the binding site, including hydrogen bonds, $\pi$-stacking, and hydrophobic contacts---consistent with known biochemistry.

\vspace{2cm}

\textbf{Scaling Behavior.} Performance improves log-linearly with training data:

\vspace{1cm}

\begin{center}
\renewcommand{\arraystretch}{1.6}
{\postertable
\begin{tabular}{@{}lc@{}}
\toprule
\textbf{Training Size} & \textbf{RMSE} \\
\midrule
1,000 complexes   & 1.72 \\
5,000 complexes   & 1.38 \\
10,000 complexes  & 1.24 \\
20,916 complexes  & 1.18 \\
\bottomrule
\end{tabular}}
\end{center}

\vspace{2cm}

% ──── Conclusions ────
\section*{Conclusions}

\posterbody

\begin{itemize}[leftmargin=1.5em, itemsep=12pt]
  \item AffinityGNN achieves \textbf{state-of-the-art} binding affinity prediction (RMSE = 1.18 kcal/mol)
  \item \textbf{Interpretable} attention mechanism reveals binding site interactions
  \item \textbf{Practical impact}: identified 4 potent M$^{\text{pro}}$ inhibitors from virtual screening
  \item Inference time: \textbf{0.3ms per complex} (GPU), enabling large-scale screening
  \item Code and models available at \texttt{github.com/oxbiochem/affinitygnn}
\end{itemize}

\vspace{2cm}

% ──── References ────
\section*{Key References}

{\fontsize{24}{32}\selectfont
\begin{enumerate}[leftmargin=1.5em, itemsep=6pt]
  \item Corso et al., ``DiffDock,'' \textit{ICLR}, 2023.
  \item St\"{a}rk et al., ``EquiBind,'' \textit{ICML}, 2022.
  \item Wang et al., ``PDBbind v2020,'' \textit{JCIM}, 2020.
  \item Kipf \& Welling, ``GCN,'' \textit{ICLR}, 2017.
  \item Veli\v{c}kovi\'{c} et al., ``GAT,'' \textit{ICLR}, 2018.
  \item Gilmer et al., ``MPNN,'' \textit{ICML}, 2017.
\end{enumerate}
}

\vspace{1.5cm}

% ──── Acknowledgements ────
\section*{Acknowledgements}

{\fontsize{24}{32}\selectfont
This work was supported by the Wellcome Trust (Grant 203141/Z/16/Z), EPSRC Doctoral Training Partnership, and computing resources from JADE2 (EP/T022205/1). We thank the Oxford Structural Genomics Consortium for providing crystal structures and binding data.}

\end{multicols}

\end{document}
